\documentclass{beamer}
\usepackage{amssymb}
\newcommand{\cinf}[1]{\text{cotinf}(#1)}

\begin{document}

  \begin{frame}
    \textit{Definición}: Sea $A \subseteq \mathbb{Q}$. Decimos que es un número real si:

    \begin{itemize}
      \item $A \neq \emptyset$.
    \end{itemize}

  \end{frame}

  \begin{frame}
    \textit{Definición}: Sea $A \subseteq \mathbb{Q}$. Decimos que es un número real si:

    \begin{itemize}
      \item $A \neq \emptyset$.
      \item $\mathbb{Q}\setminus A \neq \emptyset$.
    \end{itemize}

  \end{frame}

  \begin{frame}
    \textit{Definición}: Sea $A \subseteq \mathbb{Q}$. Decimos que es un número real si:

    \begin{itemize}
      \item $A \neq \emptyset$.
      \item $\mathbb{Q}\setminus A \neq \emptyset$.
      \item $A$ no tiene mínimo.
    \end{itemize}

  \end{frame}

  \begin{frame}
    \textit{Definición}: Sea $A \subseteq \mathbb{Q}$. Decimos que es un número real si:

    \begin{itemize}
      \item $A \neq \emptyset$.
      \item $\mathbb{Q}\setminus A \neq \emptyset$.
      \item $A$ no tiene mínimo.
      \item $\mathbb{Q}\setminus A \subseteq \cinf{A}$
    \end{itemize}

  \end{frame}

  \begin{frame}
    \textit{Definición}: Sea $A \subseteq \mathbb{Q}$. Decimos que es un número real si:

    \begin{itemize}
      \item $A \neq \emptyset$.
      \item $\mathbb{Q}\setminus A \neq \emptyset$.
      \item $A$ no tiene mínimo.
      \item $\mathbb{Q}\setminus A \subseteq \cinf{A}$
    \end{itemize}


    Denotamos el siguiente conjunto como el conjunto de los números reales.

    \[\mathbb{R} = \{A \in P(\mathbb{Q}): A \text{ es un número real}\}\]
  \end{frame}

  \begin{frame}
    
    \textit{Teorema}: Si $A$ es real, entonces $\cinf{A} = \mathbb{Q} \setminus A$.
  \end{frame}

  \begin{frame}
    \textit{Definición}: Definimos la siguiente relación sobre los números reales:
  \end{frame}

  \begin{frame}
    \textit{Definición}: Definimos la siguiente relación sobre los números reales:

    \[A \leq B \iff B \subseteq A\]
  \end{frame}

  \begin{frame}
    \textit{Definición}: Definimos la siguiente relación sobre los números reales:

    \[A \leq B \iff B \subseteq A\]

    \textit{Teorema}: La relación definida anteriormente es una relación de orden 
    total.
  \end{frame}

  \begin{frame}
    \textit{Teorema}: Si $M \subseteq \mathbb{R}$ es no vacío y acotado inferiormente, 
    tien ínfimo en $\mathbb{R}$.
  \end{frame}

  \begin{frame}
    \textit{Teorema}: Si $M \subseteq \mathbb{R}$ es no vacío y acotado inferiormente, 
    tien ínfimo en $\mathbb{R}$.
  \end{frame}

  \begin{frame}
    \textit{Teorema}: Si $M \subseteq \mathbb{R}$ es no vacío y acotado inferiormente, 
    tien ínfimo en $\mathbb{R}$.
    \newline\newline

    \textit{Definición}: Denotamos 

    \begin{equation*}
      \begin{aligned}
        0 = S^{>_\mathbb{Q}}(0_\mathbb{Q}) \\
        1 = S^{>_\mathbb{Q}}(1_\mathbb{Q})
      \end{aligned}
    \end{equation*}
  \end{frame}

  \begin{frame}
    \textit{Definición}: Definimos los conjuntos 

    \begin{itemize}
      \item $\mathbb{R}^+ = S^>(0)$.
      \item $\mathbb{R}^- = S^<(0)$.
    \end{itemize}
  \end{frame}

  \begin{frame}
    \textit{Teorema}: Si $A \in \mathbb{R}$, $A \cap \mathbb{Q} \neq \emptyset$ si y sólo si
    $A \in \mathbb{R}^-$.
  \end{frame}

  \begin{frame}
    \textit{Teorema}: Si $A \in \mathbb{R}$, $A \cap \mathbb{Q} \neq \emptyset$ si y sólo si
    $A \in \mathbb{R}^-$.
    \newline\newline

    \textit{Teorema}: Si $A \in \mathbb{R}$ y $A^c$ tiene máximo $r$, entonces 
    \[A = S^{>_\mathbb{Q}}(r)\]
  \end{frame}

  \begin{frame}
    \textit{Definición}: Sea $A \in \mathbb{R}$. Si $A^c$ tiene máximo decimos que es 
    un real racional y si no, real irracional.
  \end{frame}

  \begin{frame}
    \textit{Definición}: Sea $A \in \mathbb{R}$. Si $A^c$ tiene máximo decimos que es 
    un real racional y si no, real irracional.
    \newline\newline

    Identificamos el siguiente conjunto como los racionales reales:
    \[\mathbb{Q}_\mathbb{R} = \{S^{>_\mathbb{Q}}(r)\}_{r \in \mathbb{Q}}\]
  \end{frame}

  \begin{frame}

    \textit{Definición}: Dados $A, B \subseteq \mathbb{Q}$, definimos 
    \[A + B = \{a + b: a\in A \text{ y } b \in B\}\]

  \end{frame}

  \begin{frame}

    \textit{Definición}: Dados $A, B \subseteq \mathbb{Q}$, definimos 
    \[A + B = \{a + b: a\in A \text{ y } b \in B\}\]
    \newline\newline

    \textit{Teorema}: Dados $A$ y $B$ reales, $A + B$ es un número 
    real.

  \end{frame}

  \begin{frame}

    \textit{Definición}: Dados $A, B \subseteq \mathbb{Q}$, definimos 
    \[A + B = \{a + b: a\in A \text{ y } b \in B\}\]
    \newline\newline

    \textit{Teorema}: Dados $A$ y $B$ reales, $A + B$ es un número 
    real.
    \newline\newline

    \textit{Teorema}: para todo real, $A + 0 = A$.

  \end{frame}

  \begin{frame}
    \textit{Teorema}: para todo par de racionales $r, t$, se tiene
    \[S^{>_\mathbb{Q}}(r) +_\mathbb{R} S^{>_\mathbb{Q}}(t) = S^{>_\mathbb{Q}}(r +_\mathbb{Q} t)\]
  \end{frame}

  \begin{frame}

    \textit{Definición}: Con $A \in \mathbb{R}$, definimos

    \begin{equation*}
      -A = \left\{
      \begin{aligned}
        &S^{>_\mathbb{Q}}(-r) \text{ si }  A = S^{>_\mathbb{Q}}(r), \text{ para algún } r\in \mathbb{Q} \\
        &\{-x: x \in A^c\}, \text{ en otro caso}
      \end{aligned}
      \right.
    \end{equation*}

  \end{frame}

  \begin{frame}

    \textit{Teorema}: Sea $A \in \mathbb{R}$, $t > 0$, $c \in A^c$ y $a_0 \in A$, entonces existen 
    $b \in A^c$ y $a \in A$ tal que 

    \begin{equation*}
      \begin{aligned}
        a - b < t \qquad\qquad a_0 \ge a > b \ge c
      \end{aligned}
    \end{equation*}

  \end{frame}

  \begin{frame}

    \textit{Teorema}: Sea $A \in \mathbb{R}$, $t > 0$, $c \in A^c$ y $a_0 \in A$, entonces existen 
    $b \in A^c$ y $a \in A$ tal que 

    \begin{equation*}
      \begin{aligned}
        a - b < t \qquad\qquad a_0 \ge a > b \ge c
      \end{aligned}
    \end{equation*}
    \newline\newline

    \textit{Teorema}: Si $A \in \mathbb{R}$, $A + (-A) = 0$.
  \end{frame}

  \begin{frame}

    \textit{Teorema}: Sean $A, B \in \mathbb{R}$, entonces 
    \[A \leq B \iff B + (-A) \in \mathbb{R}^+\]
  \end{frame}

\end{document}
